\section{Result III: The Audit}
\label{sec:audit}
% Every percentage: "a purposive sample of 20", never "the field". No paper named as a failure.

\subsection{Method}
\label{sec:audit-method}
We scored a purposive sample of 20 papers (2018--2026) that make a comparative architectural claim
about a learned performance or cost model on five rigor axes: (1)~seed variance, (2)~parameter
matching, (3)~significance testing, (4)~evaluation-set integrity, and (5)~effect size against noise and
baseline strength. Each axis is scored pass, partial or fail, and every score is backed by a quoted
passage.

Papers were selected by snowball sampling from the lineage of learned cost model papers plus searches;
the search terms were not logged and the selection protocol was written afterwards. All 20 papers were
first scored by AI. Five were then human-verified on all five axes, which corrected 3 of 25 scores, all
towards stricter; Axis 3 was human-verified for all 20 because the pre-screen failed every paper on it.
The only score that became looser was TPUGraphs, from fail to partial, because it reports $\pm1\sigma$
but runs no test. For the other 15 papers, Axes 1, 2, 4 and 5 are AI-prescreened only. To estimate
coverage, 6 fixed Semantic Scholar queries returned 571 candidates; our sample includes 8 of the 30
most-cited in-scope papers and 0 of the 10 most-cited from 2024--2026. That screening retrieves only 10
of our 20 papers (it misses TPUGraphs), and in-scope papers outside the sample were not audited.

\subsection{No paper runs a significance test}
\label{sec:audit-result}
None of the 20 papers runs a significance test on its headline architectural claim, 17 (85\%) fail at
least three of the five axes, and none passes Axes 1--3 together (Table~\ref{tab:audit}).
% Frame B (O15): sentence fixed in advance by the frozen declaration (section 6, "0 of n"), counts 0/10.
In a separate, pre-declared check of the 10 most-cited in-scope papers from 2024--2026 in the same
frame, none of the 10 we could score runs a significance test for its headline comparison
(Appendix~D).

\begin{table}[t]
\caption{Audit outcome by rigor axis for 20 papers that make a comparative architectural claim about a
learned performance or cost model (2018--2026). No paper passes Axes 1--3 together, and 17 of 20 fail
at least three of the five axes.}
\label{tab:audit}
\centering
\small
% Generated by scripts/paper_figures.py from the per-paper tables in experiments/phase7_meta_analysis.md
% papers failing >=3 axes: 18/20; passing axes 1-3: 0/20
\begin{tabular}{lccc}
\toprule
Axis & Fail & Partial & Pass \\
\midrule
1. Seeds / variance & 16/20 & 4/20 & 0/20 \\
2. Parameter-matched & 17/20 & 3/20 & 0/20 \\
3. Significance test & 20/20 & 0/20 & 0/20 \\
4. Eval-set integrity & 4/20 & 7/20 & 9/20 \\
5. Effect vs noise & 6/20 & 14/20 & 0/20 \\
\bottomrule
\end{tabular}

\end{table}

\subsection{Typical practice cannot resolve effects of this size}
\label{sec:audit-resolve}
These failures are not missing paperwork. No paper reports the seed-to-seed spread of its headline
comparison, so readers cannot discern whether a gain exceeds the noise. By Table~\ref{tab:mde}, five seeds
resolve only $\Delta\tau \approx$ 0.07--0.12, and the effects in Sections~\ref{sec:result1}
and~\ref{sec:capstone} are smaller. Our own results show both failure modes. The DiffPool win only came from
capacity and vanished at matched size (Section~\ref{sec:protocol}). The $+0.165$ contention-edge result
had five seeds, matched parameters and $p < .001$, passing Axes 1--3, yet only the pre-registered
baseline check exposed it as a learning-rate handicap (Section~\ref{sec:tuning}). Passing every check
the field currently applies is not enough.

\subsection{Good practice already exists}
\label{sec:audit-good}
The problem is a habit, not a rule. Kaufman et al.\ report a case where their own learned model loses to
the analytical baseline on the harder split of their tile-size task (6.3\% against 2.3\% mean
error)~\cite{kaufman2021learned}. PerfVec declines a head-to-head state-of-the-art claim because the
hardware differed~\cite{li2024learning}. HLS-GNN runs a five-seed protocol~\cite{wu2022highlevel}. Rigorous comparison is achievable within current practice.
